\begin{itemize}
    \item[а)] $L$ - поле
    \par \Proof Рассмотрим $a,b\neq 0 \in L$. Тогда $\Q(a, b)$ — алгебраическое расширение $\Rightarrow a+b, a \cdot b, \frac{a}{b} \in L \Rightarrow L$ — поле. \EndProof
    
    \item[б)] $L$ - алгебраически замкнутое поле
    \par \Proof Рассмотрим $f(x) \in L[x], \deg f \geq 1$. Он раскладывается на линейные множители в $\Comp$. Тогда $L(\alpha) \supset L$ — алгебраическое расширение, поэтому $L(\alpha) \supset \Q$ — алгебраическое расширение, значит $\alpha$ алгебраичен над $\Q$, значит $\alpha \in L$. Значит все корни $f(x)$ лежат в $L \Rightarrow L$ алгебраически замкнуто. \EndProof
    
    \item[в)] $L=\overline{\Q}$
    \par \Proof Так как расширение $L \supset \Q$ алгебраическое и $L$ алгебраически замкнуто, по определению $L=\overline{\Q}$ \EndProof
\end{itemize}